%%%%%%%%%%%%%%%%%%%%%%%%%%%%%%%%%%%%%%%%%%%%%%%%%%%%%%%%%%%%%%%%%%%%%%%%%%%
%
% Generic template for TFC/TFM/TFG/Tesis
%
% $Id: introduccion.tex,v 1.1 2015/06/05 00:16:17 macias Exp $
%
% By:
%  + Javier Macías-Guarasa.
%    Departamento de Electrónica
%    Universidad de Alcalá
%  + Roberto Barra-Chicote.
%    Departamento de Ingeniería Electrónica
%    Universidad Politécnica de Madrid
%
% Based on original sources by Roberto Barra, Manuel Ocaña, Jesús Nuevo,
% Pedro Revenga, Fernando Herránz and Noelia Hernández. Thanks a lot to
% all of them, and to the many anonymous contributors found (thanks to
% google) that provided help in setting all this up.
%
% See also the additionalContributors.txt file to check the name of
% additional contributors to this work.
%
% If you think you can add pieces of relevant/useful examples,
% improvements, please contact us at (macias@depeca.uah.es)
%
% You can freely use this template and please contribute with
% comments or suggestions!!!
%
%%%%%%%%%%%%%%%%%%%%%%%%%%%%%%%%%%%%%%%%%%%%%%%%%%%%%%%%%%%%%%%%%%%%%%%%%%%

\chapter{Introducción}
\label{cha:introduction}

\begin{FraseCelebre}
  \begin{Frase}
    Desocupado lector, sin juramento me podrás creer que quisiera que este
    libro [...] fuera el más hermoso, el más gallardo y más discreto que
    pudiera imaginarse\footnote{Tomado de ejemplos del proyecto \texis{}.}.
  \end{Frase}
  \begin{Fuente}
    Miguel de Cervantes, Don Quijote de la Mancha
  \end{Fuente}
\end{FraseCelebre}


\section{Presentación}
\label{sec:intro-presentacion}

Este capítulo contiene algunos ejemplos básicos de los elementos que pueden
utilizarse en la elaboración de la memoria.

La figura~\ref{fig:intro-figura} muestra un ejemplo de inclusión de una figura.

\begin{figure}[h]
  \centering
  \includegraphics[width=0.55\textwidth]{Figure1}
  \caption{Ejemplo de figura convencional.}
  \label{fig:intro-figura}
\end{figure}

La tabla~\ref{tab:intro-tabla} muestra un ejemplo de una tabla convencional.

\begin{table}
  \renewcommand{\arraystretch}{1.3}
  \caption{Comparativa.}
  \label{tab:intro-tabla}
  \begin{center}
    \begin{tabular}{|c|c|c|}
      \hline
      Method & Training Time & Man-Work (\%)\\
      \hline
      Propagation model & $<$ 30 sec & 5\\
      \hline
      Manual & 9 h 30 min & 24\\
      \hline
      Automatic & 2 h & 10\\
      \hline
    \end{tabular}
  \end{center}
\end{table}

También pueden incluirse expresiones matemáticas numeradas, como la
ecuación~\ref{eq:intro-ejemplo}:

\begin{equation}
  \label{eq:intro-ejemplo}
  p[q_t=\sigma_t\mid q_{t-1}=\sigma_{t-1}]
\end{equation}

Finalmente, la plantilla permite gestionar automáticamente acrónimos. Por
ejemplo, \ac{ETTS} está definido en el fichero de acrónimos de la plantilla y aparece la primera vez en su versión expandida seguida del acrónimo entre paréntesis, y cuando lo utilizas una segunda vez (como aquí \ac{ETTS}), sólo verás el acrónimo.


\section{Organización de la memoria}
\label{sec:intro-organizacion-memoria}

Esta memoria se organiza en \ldots grandes capítulos. El primero \ldots


%%% Local Variables:
%%% TeX-master: "../book"
%%% End:
